\input{template}

\title{\bfseries 基于 Dijkstra 算法的地铁网络最短路径规划}
\author{韦晏洲 \quad PB24000379}

\begin{document}
	\maketitle
	\begin{abstract}
		本文针对城市地铁出行中的路径规划问题，实现了一个基于 Dijkstra 算法的交互式最短路径检索系统。为了准确模拟换乘行为，模型引入了“拆点”技术，将每个物理站点根据所属线路拆分为独立的逻辑节点，并建立等效时间权重的换乘边，从而将带换乘惩罚的路径搜索转化为经典图论中的最短路问题。实验测试表明，该模型能够有效处理大规模线网数据，准确规避高成本换乘，并输出具备实际参考价值的导航指令。最后，本文对图布局失真、引入实时列车时刻表的时空动态规划算法等改进方向进行了探讨。
		
		\textit{关键词：地铁路径规划；Dijkstra 算法；拆点模型；换乘优化；可视化系统}
	\end{abstract}
	\clearpage
	\section{前言}
	
	在城市轨道交通系统中，路径规划是导航系统与运营调度的核心问题。本报告旨在针对地铁线网模型，通过数学建模与算法实现，探索两站之间的最短路径求解方案。本项任务分为两个阶段进行探讨。
	
	\subsection{基础任务}
	
	我们将地铁网络抽象为标准的无向加权图 $G = (V, E)$。在此模型下：
	
	\begin{itemize}
	\item 节点 $V$：代表物理上的地铁站点。

	\item 边 $E$：代表站点间的轨道连接，其权重给定。
	\end{itemize}

	该任务的核心目标是实现 Dijkstra 算法，验证算法在无向加权图中的准确性。此模型虽然能够处理点对点的连通性问题，但在处理换乘问题时存在局限。
	
	\subsection{进阶任务}
	
	进阶任务引入了换乘时间这一变量。由于基础模型中一个节点仅代表一个物理位置，无法直接在单一节点上表达“换乘”。为此，我们引入了“拆点”技术：将单一换乘站拆分为对应不同线路的多个逻辑节点。在这些逻辑节点之间建立虚拟换乘边，并将换乘时间赋权给这些边，从而使 Dijkstra 算法能够在不改变核心逻辑的前提下，求解出包含换乘时间在内的全局最优路径。
	
	\section{相关工作}
	Dijkstra 算法由荷兰计算机科学家 Edsger W. Dijkstra 于 1956 年提出，并于 1959 年正式发表，是图论中最经典的单源最短路径算法之一。该算法在计算机网络、地图导航以及各类运筹优化问题中有着极其广泛的应用。
	
	\section{建立数学模型}
	
	\subsection{基础任务}
	
	在不考虑换乘时间的情况下，地铁网络可抽象为一个无向加权图 $G = (V, E, C)$。在该模型中，我们假设乘客在换乘站改变线路时不产生任何额外的时间成本，仅考虑站点间的区间运行耗时。将地铁网络转化为如下的图：
	
	\begin{itemize}
	\item 顶点集 $V = \{1, 2, \dots, N\}$：其中每一个整数 $i \in V$ 唯一对应一个物理地铁站点，共有 $N$ 个站点。
	
	\item 边集 $E \subseteq \{ \{i, j\} \mid i, j \in V, i \neq j \}$：若站点 $i$ 与站点 $j$ 之间存在直接的轨道连接，则定义无向边 $\{i, j\} \in E$ \footnote{Dijkstra算法允许有向边，\{$i \to j$\} 和 \{$j \to i$\}的赋值也允许不相同，不过在此任务中无需考虑。}。
	
	\item 权重集合 $C$：对每一条边 $\{i, j\} \in E$，分配一个正权值 $c_{ij} \in \mathbb{R}^+$ \footnote{Dijkstra算法允许权值为0（但不允许为负），不过在此任务中无需考虑。}，表示列车在站点 $i$ 与 $j$ 之间的运行耗时。由于图 $G$ 是无向的，权值满足对称性，即 $c_{ij} = c_{ji}$
	\end{itemize}
	
	则该问题转化为：
	
	给定起点 $s \in V$ 与终点 $t \in V$，寻找一条路径 $P = (i_1, i_2 \cdots i_k)$，满足$i_1 = s, i_k = t$，且最小化路径上所有边的权值之和，即$$ \boxed{P = \argmin_{P} \sum_{n=1}^{k-1} c_{i_{n-1}, i_n}}$$
	
	同时为了表述方便，将这个最小值成为 $s$ 与 $t$ 之间的距离 $\dist(s,t)$.
	
	\subsection{进阶任务}
	
	为了考虑换乘时间成本，基础模型中的单一节点已无法表示换成车站。因此，我们构建\textit{扩展图} $G' = (V', E', C')$，通过“拆点”技术进行建模。
	
	\begin{itemize}
		\item {顶点集 $V'$}：设经过物理站点 $i$ 的线路集合为 $L(i)$。我们将每个物理站点 $i \in V$ 拆分为多个逻辑顶点：
		$$V' = \{ i^{(l)} \mid i \in V, \ l \in L(i) \}$$
		其中 $i^{(l)}$ 表示“位于线路 $l$ 上的站点 $i$”。
		
		\item {边集 $E'$} 由以下两部分组成：
		\begin{enumerate}
			\item {运行边}：若物理站点 $i, j$ 在线路 $l$ 上相邻，则存在边 $\{i^{(l)}, j^{(l)}\} \in E'$，其权重为 $c'_{i^{(l)}, j^{(l)}} = c_{ij}$。
			\item {换乘边}：对于任意物理站点 $i$，若其属于两条不同线路 $l$ 与 $r$ ($l \neq r$)，则在此基础上添加虚拟边 $\{i^{(l)}, i^{(r)}\} \in E'$，其权重 $c'_{i^{(l)}, i^{(r)}} = T_{transfer}$，代表换乘耗时。
		\end{enumerate}
	\end{itemize}
	
	在该模型下，给定起点物理站 $s$ 与终点物理站 $t$，目标转化为在 $G'$ 中寻找连接集合 $S = \{s^{(l)} \mid l \in L(s)\}$ 与集合 $T = \{t^{(r)} \mid r \in L(t)\}$ 的全局最短路径：
	$$ P = \argmin_P \ \{\dist(u, v): u \in S, v \in T, P \text{ 连接$u$与$v$}\} $$
	
	\paragraph{换乘站权重设计}
	在进阶任务中我们需要考虑换乘站带来的时间并最小化路径总时长。由于地铁旅速未知，我们无法得到相邻站点之间确切的运行时长，所以只可利用站距/旅速进行估算。资料显示，市域普通线路平均旅速大致为 $35km/h$ \footnote{见B站视频 \href{https://www.bilibili.com/video/BV179jqzfE18}{\textit{北京地铁平均旅速痛度表}}，我们这里不再区分各条线路之间的旅速，因为（2009年前开通的）各条线路旅速都在 $30-40km/h$，且估算换乘时间也带来较大误差。}。换乘步行、等车时间各按2.5分钟计算，因此换乘时长估计为5分钟，换乘站之间的权重按$2.917km$计算。
	
	

	\section{Dijkstra算法描述}
	
	\begin{algorithm}[H]
		\caption{Dijkstra 算法}
		\begin{algorithmic}[1]
			\State \textbf{输入：} 图 $G = (V, E)$（无向或有向），非负边权 $c: E \to \mathbb{R}_{\geq 0}$，源点 $s \in V$。
			\State \textbf{输出：} 所有顶点的最短路长度 $\dist(s, v)$。
			\State \textbf{初始化：} $d(s) \gets 0$；对 $v \neq s$ 令 $d(v) \gets \infty$；已确定集合 $X \gets \emptyset$。
			\While{$X \neq V$}
			\State 取 $v^* \in V \setminus X$ 使 $d(v^*) = \min\{d(v) : v \in V \setminus X\}$。
			\State 令 $X \leftarrow X \cup \{v^*\}$ \quad \{此时 $d(v^*) = \dist(s, v^*)$，见命题~\ref{prop1}\}
			\For{每条边 $(v^*, w) \in E$}
			\If{$d(w) > d(v^*) + c(v^*, w)$}
			\State $d(w) \gets d(v^*) + c(v^*, w)$
			\EndIf
			\EndFor
			\EndWhile
		\end{algorithmic}
	\end{algorithm}
	
	\begin{prop*}{}{prop1}
		在执行上述算法时，恒有 $\forall v \in X, d(v) = \dist(v,x)$.
	\end{prop*}
	
	\begin{pf}
	假设算法执行到某一轮，第 5 行选出的 $v^*$ \textit{第一次}未达到真实最短距离，即 $\dist(s, v^*) < d(v^*)$。此时，考虑连接 $s$ 到 $v^*$ 的一条真实最短路径，取该路径上第一条离开集合 $X$ 的边 $(y', y)$（其中 $y' \in X, y \notin X$）。
	
	由于最短路径的子路径依然是最短路径，根据最优性原理有：
	\[ \dist(s, y) = \dist(s, y') + c(y', y) = d(y') + c(y', y) \]
	因为 $y' \in X$，根据算法的更新规则，节点 $y$ 的当前估算距离满足 $d(y) \leq d(y') + c(y', y) = \dist(s, y)$。结合最短路径的定义可知 $d(y) = \dist(s, y)$。
	
	又因为在算法第 5 行中，我们选择了所有未确定节点中 $d(\cdot)$ 最小的 $v^*$，因此有：
	\[ d(v^*) \leq d(y) = \dist(s, y) \leq \dist(s, v^*) < d(v^*) \]
	其中，$\dist(s, y) \leq \dist(s, v^*)$ 是由于边权 $c$ 的非负性保证了最短路径长度随节点增加而单调不减。
	
	由此得到 $d(v^*) < d(v^*)$，的矛盾。故 $d(v^*) = \dist(s, v^*)$。因此，将 $v^*$ 加入集合 $X$ 后，命题依然成立。
	\end{pf}
	
	\section{实验结果}
	
	通过互联网考古，我们找到了2009年的北京地铁图、伦敦地铁图和2010年的上海地铁图。
	\fig{1}{1}{2009年北京地铁图}
	\fig{2}{1}{2010年上海地铁图}
	\fig{3}{1}{2009年伦敦地铁图}
	
	\subsection{基础任务}
	我们对北京、上海、伦敦三个城市中距离较远的站点分别进行了测试，通过比对各自地铁图，通过肉眼观察，程序输出的最短路径无问题。详细代码见\texttt{Code}文件夹。
	
	在对北京地铁的测试中，我们发现从4号线安河桥北站到八通线八里桥站路径进行了多次换乘（4-10-13-2-1-八通），这在实际搭乘过程中并不现实。因此，我们需要对换乘施加惩罚项。
	
	\fig{test1}{1}{北京地铁测试样例}
	\fig{test2}{1}{上海地铁测试样例}
	\fig{test3}{1}{伦敦地铁测试样例}
	
	\subsection{拓展任务}
	我们依然对北京地铁中距离较远的站点进行测试。为了优化体验，我们在路径输出时，在换乘站标明了“Interchage to Line X”.在测试一中，我们同样测试了四号线安河桥北到八通线八里桥的最短路径，可以发现此时算法给出的最有路径确实换乘最少。详细代码见 \texttt{Code Optional} 文件夹。
	
	\fig{test'1}{1}{安河桥北--八里桥}
	\fig{test'2}{1}{西直门--芍药居}
	\fig{test'3}{1}{天通苑北--公益西桥}
	
	
	\section{改进方向与策略}
	
	\subsection{Fruchterman-Reingold 弹簧算法的不足}
	在本实验的 GUI 实现中，采用了Fruchterman-Reingold 弹簧算法来生成站点的可视化坐标。虽然该算法能够清晰地展示图 $G=(V, E)$ 的拓扑连通性，但在处理实际地铁线网时与各站点的这很是地理位置存在严重不符。
	
	同时，由于布局算法依赖随机种子，同一线网在不同次运行中可能呈现完全不同的形状，缺乏确定性。
	
	\paragraph{改进策略}
	一个经典的方法是在原始数据集中增加经纬度信息 $(\text{lon, lat})$：将经度与纬度转换为界面坐标系的 $(x, y)$。$$x_i = \frac{\text{lon}_i - \text{lon}_{\min}}{\text{lon}_{\max} - \text{lon}_{\min}}, \quad y_i = \frac{\text{lat}_i - \text{lat}_{\min}}{\text{lat}_{\max} - \text{lat}_{\min}}$$
	
	当然也可以把地铁图导入界面。一个比较好的用户界面是\href{https://www.sz-mtr.com/service/guide/map/index_zh.html?T}{\textit{苏州地铁--时长图}}。
	
	\subsection{换乘时间数据的不足}
	
	在建立换乘图模型时，由于原始数据并未包含各换乘站内部的实际步行距离及垂直电梯转换时间，模型不得不采取了简化假设，将所有换乘边的权重统一设定为固定常数 $\Delta t = 5$ 分钟。这一设定忽视了站点结构的不同：在实际地铁网络中，换乘时间受车站建筑结构影响极大。例如，同台换乘可能仅需 30 秒，而像北京平安里、望京西等站点，不同线路站台间的距离可能导致换乘耗时超过 8--10 分钟。统一设定为 5 分钟忽略了这种极差，可能导致算法在推荐路径时产生偏差。
	
	如果未来在这一方向改进，我们将需要更多相关数据。
	
	\subsection{动态规划}
	在网站\href{https://bjsubway.boyinthesun.cn/}{\textit{北京地铁可视化}}中，可以查询到北京地铁的列车时刻表。如果在知道各个站点换乘时间的基础上引入所有列车的时刻表，将可以实现分钟级导航。
	
	一个经典的动态规划算法（Connection Scan Algorithm (CSA)）如下。其核心思路是将时刻表的每一趟列车按相邻两站的出发与达到拆为“元列车”（如 19:23 天安门东 --- 19:25 天安门西），然后从计划出发时间开始，按“元列车”出发的时间从早到晚搜索。如果该“元列车”可以对到达时间做出改进，则更新记录。直至搜索完 $EarliestArr$ 之前的全部列车。
	
	\begin{algorithm}[H]
		\caption{基于列车时刻表的动态规划最短路径算法}
		\begin{algorithmic}[1]
			\Require 站点集合 $V$, 列车班次集合 $C$, 起点 $s$, 终点 $t$, 出发时间 $T_{start}$
			\Ensure 到达终点的最早时间 $EarliestArr[t]$
			
			\State \textbf{Initialization:} 
			\State 令 $EarliestArr[s] = T_{start}$，其余所有站点 $EarliestArr[v] = \infty$
			\State 将所有班次 $c \in C$ 按发车时间 $T_{dep}(c)$ 进行升序排列
			
			\State \textbf{DP State Transition (Connection Scanning):}
			\For{每个班次 $c \in C$ (由车站 $u$ 开往车站 $v$)}
			\State $T_{dep} \gets$ 班次 $c$ 从 $u$ 站的发车时间
			\State $T_{arr} \gets$ 班次 $c$ 到达 $v$ 站的到达时间
			
			\If{$EarliestArr[u] + T_{transfer}(u) \leq T_{dep}$} \Comment{判断是否能赶上这班车}
			\If{$T_{arr} < EarliestArr[v]$}
			\State $EarliestArr[v] \gets T_{arr}$ \Comment{更新到达 $v$ 站的最早时间}
			\State $Parent[v] \gets c$ \Comment{记录捕获的班次用于路径回溯}
			\EndIf
			\EndIf
			\EndFor
			
			\State \Return $EarliestArr[t]$
		\end{algorithmic}
	\end{algorithm}
	
	不过其实施较为复杂，故不在本实验中完成。
	
	\paragraph{真实的导航是怎么处理的？}
	
	工业级的导航系统（如高德、Google Maps）能实现秒级响应，除了上述Dijkstra、动态规划算法，还有下面一些重要算法：\footnote{这一段只是我单纯好奇，结合Gemini进行了一个简单调研，与实验报告关系不大了。}
	
	\begin{enumerate}
		\item 空间分层：等级化搜索。
		
		算法不会在搜索时考虑所有的路。它像人类的大脑一样有“远近观念”：
		
		\textit{长距离}：如果你从上海去北京，算法在初始阶段只看“高速公路”和“国道”这种骨干网，完全忽略居民区里的小路。
		
		\textit{短距离}：只有当你接近目的地时，算法才会把搜索精度降到“街道”级别。
		
		通过这种方式，搜索范围从“一整张密密麻麻的网”缩减成了“几条高速 + 两端的局部路网”。
		
		\item 空间分区：多层 Dijkstra。
		
		工程师会将整个地图切成很多个“区域”。
		
		\textit{预计算}：算法会提前算好每个区域“边界点”到“边界点”的最短距离。
		
		\textit{跨区搜索}：当你跨区旅行时，算法不再进入区域内部数星星，而是直接像跳格子一样在边界点之间跳转。
		
		这就像去另一个城市旅游，只需要知道怎么到火车站，以及火车站之间怎么坐车，而不需要知道中间经过的每个小镇该怎么走。
		
		\item 针对公交、地铁的RAPTOR 算法。
		
		对于上文中的时刻表动态规划，在工业界有一个著名的改进版叫 RAPTOR (Round-Based Public Transit Routing)。它和 Dijkstra 完全不同：
		
		\textit{按“轮”搜索}：第一轮找坐 1 次车能到的站，第二轮找换乘 1 次能到的站……
		
		\textit{无视边权}：它不看具体的距离，只看线路。因为它发现：地铁和公交的换乘次数通常不会超过 4-5 次。
		
		这种算法不需要维护复杂的“优先队列”，搜索速度极快，是目前处理海量公交、地铁数据的标配。
		
		\item 预计算：收缩层级。
		
		这是目前最极端的优化手段。
		系统会提前将地图中那些“不重要”的节点（比如只有两条路相连的过路点）“收缩”掉，并给跳过的路径加上一条虚拟的“快捷边”。
		最终在线查询时，算法面对的是一个被极度简化的、全都是“快捷边”的图，从而加快搜索速度。
	\end{enumerate}
	
	
	\subsection{数据集的小错误}
	
	在完成实验时，我们注意到数据集\texttt{Beijing-2010-station-lines.txt}存在小错误。部分站点出现与未开通的6,7,9,14,15号线换乘，但是在\texttt{Beijing-2010-adjacency-distance.csv}中并未记录这些未开通线路的链接。该错误会增加扩展图的节点个数，但不会影响实验的完成。
	

\end{document}


